\Needspace{8\baselineskip}
\Article{273}
Here are some examples to illustrate the preceding rules more fully.

\emph{Example} 1. Let \(f=\left(\begin{smallmatrix}19,&21,&50\\15,&28,&1\end{smallmatrix}\right)\); then \(F=\left(\begin{smallmatrix}-825,&-166,&-398\\257,&573,&-370\end{smallmatrix}\right)\), \(D=-1\). Since \((19,1,21)\) is a reduced binary form to which no other form with a first coefficient smaller than \(19\) is equivalent, the first reduction is not applicable here. But the theory of equivalence of binary forms shows that the binary form
\SourceLeaf{316}
\((A'',B,A')=(-398,257,-166)\) can be transformed into the simpler equivalent form \((-2,1,-10)\), into which it passes by the substitution \(2,7,3,11\). So, setting \(\beta'=2\), \(\gamma'=-7\), \(\beta''=-3\), \(\gamma''=11\), we must apply to the form \(f\) the substitution \(\left\{\begin{smallmatrix}1,&0,&0\\0,&2,&-7\\0,&-3,&11\end{smallmatrix}\right\}\), by which it is found to pass into \(\left(\begin{smallmatrix}19,&354,&4769\\-1299,&301,&-82\end{smallmatrix}\right)\ldots f'\). The third coefficient of the adjoint of this form is \(-2\), so in this respect \(f'\) must be regarded as simpler than \(f\).

The first reduction can be applied to the form \(f'\). For since the binary form \((19,-82,354)\) is transformed into \((1,0,2)\) by the substitution \(13,4,3,1\): we must apply to the form \(f'\) the substitution \(\left\{\begin{smallmatrix}13,&4,&0\\3,&1,&0\\0,&0,&1\end{smallmatrix}\right\}\), by which it passes into \(\left(\begin{smallmatrix}1,&2,&4769\\-95,&16,&0\end{smallmatrix}\right)\ldots f''\).

The second reduction can again be applied to the form \(f''\), whose adjoint is \(\left(\begin{smallmatrix}-513,&-4513,&-2\\-95,&32,&1520\end{smallmatrix}\right)\). For \((-2,-95,-4513)\) passes by the substitution \(47,1,-1,0\) into \((-1,1,-2)\): so we must apply to \(f''\) the substitution \(\left\{\begin{smallmatrix}1,&0,&0\\0,&47,&-1\\0,&1,&0\end{smallmatrix}\right\}\), by which it passes into \(\left(\begin{smallmatrix}1,&257,&2\\1,&0,&16\end{smallmatrix}\right)\ldots f'''\). Its first coefficient cannot be decreased any further by the first reduction, nor can the third coefficient of its adjoint by the second.

\emph{Example} 2. Let the given form be \(\left(\begin{smallmatrix}10,&26,&2\\7,&0,&4\end{smallmatrix}\right)\ldots f\), whose adjoint is \(\left(\begin{smallmatrix}-3,&-20,&-244\\70,&-28,&8\end{smallmatrix}\right)\) and whose determinant \(=2\). Applying the second and first reductions alternately gives, in succession,
\[
\begin{array}{c|c|c}
\text{substitutions}&\text{take}&\text{into}\\[3pt]
\left\{\begin{smallmatrix}1,&0,&0\\0,&-1,&0\\0,&4,&-1\end{smallmatrix}\right\}
&f&\left(\begin{smallmatrix}10,&2,&2\\-1,&0,&-4\end{smallmatrix}\right)=f'\\[6pt]
\left\{\begin{smallmatrix}0,&-1,&0\\1,&-2,&0\\0,&0,&1\end{smallmatrix}\right\}
&f'&\left(\begin{smallmatrix}2,&2,&2\\2,&-1,&0\end{smallmatrix}\right)=f''\\[6pt]
\left\{\begin{smallmatrix}1,&0,&0\\0,&-1,&0\\0,&2,&-1\end{smallmatrix}\right\}
&f''&\left(\begin{smallmatrix}2,&2,&2\\-2,&1,&-2\end{smallmatrix}\right)=f'''\\[6pt]
\left\{\begin{smallmatrix}1,&0,&0\\1,&1,&0\\0,&0,&1\end{smallmatrix}\right\}
&f'''&\left(\begin{smallmatrix}0,&2,&2\\-2,&-1,&0\end{smallmatrix}\right)=f''''
\end{array}
\]
The form \(f''''\) cannot be reduced any further by the first or second reduction.

\Needspace{8\baselineskip}
\Article{274}
When we have a ternary form whose first coefficient, and the third coefficient of its adjoint, have been lowered as far as possible by the preceding methods: the following method gives a further reduction.

\SourceLeaf{317}
Using the same notation as in Article 272, and setting \(\alpha=1\), \(\alpha'=0\), \(\beta'=1\), \(\alpha''=0\), \(\beta''=0\), \(\gamma''=1\), \emph{that is}, using the substitution
\[
\begin{array}{ccc}
1,&\beta,&\gamma\\
0,&1,&\gamma'\\
0,&0,&1
\end{array}
\]
we have
\[
\begin{aligned}
m&=a,\quad m'=a'+2b''\beta+a\beta\beta,\\
m''&=a''+2b\gamma'+2b'\gamma+a\gamma\gamma+2b''\gamma\gamma'+a'\gamma'\gamma'\\
n&=b+a'\gamma'+b'\beta+b''(\gamma+\beta\gamma')+a\beta\gamma,\\
n'&=b'+a\gamma+b''\gamma',\quad n''=b''+a\beta
\end{aligned}
\]
and also
\[
M''=A'',\quad N=B-A''\gamma',\quad N'=B'-N\beta-A''\gamma
\]
Thus a substitution of this kind leaves unchanged the coefficients \(a\), \(A''\) which were decreased by the preceding reductions. Our task is therefore to obtain decreases in the remaining coefficients by a suitable choice of \(\beta\), \(\gamma\), \(\gamma'\). For this purpose we first observe that if \(A''=0\), we may assume that \(a=0\) as well. For if \(a\) is not \(=0\), the first reduction could still be applied once more, since every binary form of determinant \(0\) is equivalent to a form such as \((0,0,h)\), or one whose first coefficient \(=0\) (see Article 215). In exactly the same way we may assume that \(A''=0\) as well if \(a=0\), so that either neither of the numbers \(a\), \(A''\) is \(0\), or both are.

In the first case it is clear that \(\beta\), \(\gamma\), \(\gamma'\) can be chosen so that, without regard to sign, \(n''\), \(N\), \(N'\) are no greater, respectively, than \(\frac{1}{2}a\), \(\frac{1}{2}A''\), \(\frac{1}{2}A''\). Thus in the first example of the preceding article, the last form \(\left(\begin{smallmatrix}1,&257,&2\\1,&0,&16\end{smallmatrix}\right)\), whose adjoint is \(\left(\begin{smallmatrix}-513,&-2,&-1\\1,&-16,&32\end{smallmatrix}\right)\), passes by the substitution \(\left\{\begin{smallmatrix}1,&-16,&16\\0,&1,&-1\\0,&0,&1\end{smallmatrix}\right\}\) into \(\left(\begin{smallmatrix}1,&1,&1\\0,&0,&0\end{smallmatrix}\right)\ldots f''''\), whose adjoint is \(\left(\begin{smallmatrix}-1,&-1,&-1\\0,&0,&0\end{smallmatrix}\right)\).

In the second case, where \(a=A''=0\), and therefore also \(b''=0\), we have
\[
\begin{aligned}
m&=0,\quad m'=a',\quad m''=a''+2b\gamma'+2b'\gamma+a'\gamma'\gamma'\\
n&=b+a'\gamma'+b'\beta,\quad n'=b',\quad n''=0
\end{aligned}
\]
So we will have
\[
D=a'b'b'=m'n'n'
\]
and it is easy to see that \(\beta\) and \(\gamma'\) can be chosen so that \(n\) equals the residue of \(b\) of smallest absolute value modulo the greatest common divisor
\SourceLeaf{318}
of \(a'\), \(b'\); that is, so that \(n\) is no greater than half this divisor without regard to sign, and hence \(n=0\) whenever \(a'\), \(b'\) are coprime. Having determined \(\beta\), \(\gamma'\) in this way, we can choose the value of \(\gamma\) so that \(m''\) is no greater than \(b'\), without regard to sign. This would indeed be impossible if \(b'=0\), but then we would have \(D=0\), a case we have excluded. Thus for the last form in \emph{Example} 2 of the preceding article we obtain \(n=-2-\beta+2\gamma'\), so setting \(\beta=-2\), \(\gamma'=0\) gives \(n=0\). Further, \(m''=2-2\gamma\), and setting \(\gamma=1\) gives \(m''=0\). We therefore have the substitution \(\left\{\begin{smallmatrix}1,&-2,&1\\0,&1,&0\\0,&0,&1\end{smallmatrix}\right\}\), by which that form passes into \(\left(\begin{smallmatrix}0,&2,&0\\0,&-1,&0\end{smallmatrix}\right)\ldots f'''''\).

\Needspace{8\baselineskip}
\Article{275}
Suppose we have a sequence of equivalent ternary forms \(f\), \(f'\), \(f''\), \(f'''\), etc., and transformations of each of these forms into the next. From the transformations of the form \(f\) into \(f'\), and of the form \(f'\) into \(f''\), Article 270 gives a transformation of the form \(f\) into \(f''\). From this and the transformation of the form \(f''\) into \(f'''\) follows a transformation of the form \(f\) into \(f'''\), etc.\ Clearly, in this way a transformation of the form \(f\) into any other in the sequence can be found. And since a transformation of the form \(f\) into any other equivalent form \(g\) gives a transformation of the form \(g\) into \(f\) (\(S''\) from \(S\), Articles 268, 269), we can in this way obtain a transformation of any form in the sequence \(f'\), \(f''\), etc.\ into the first, \(f\). --- Thus for the forms in the first example of the preceding article we find the substitutions
\[
\begin{array}{rrr|rrr|rrr}
13,&4,&0&13,&188,&-4&13,&-20,&16\\
6,&2,&-7&6,&87,&-2&6,&-9,&7\\
-9,&-3,&11&-9,&-130,&3&-9,&14,&-11
\end{array}
\]
by which \(f\) passes into \(f''\), \(f'''\), \(f''''\), respectively, and from the last substitution we obtain \(\left\{\begin{smallmatrix}1,&4,&4\\3,&1,&5\\3,&-2,&3\end{smallmatrix}\right\}\), by which \(f''''\) passes into \(f\). Similarly, for \emph{Example} 2 of the preceding article, we obtain the substitutions
\[
\begin{array}{rrr|rrr}
1,&-1,&1&2,&-3,&-1\\
-3,&4,&-3&3,&1,&0\\
10,&-14,&11&2,&4,&1
\end{array}
\]
by which, respectively, the form \(\left(\begin{smallmatrix}10,&26,&2\\7,&0,&4\end{smallmatrix}\right)\) passes into \(\left(\begin{smallmatrix}0,&2,&0\\0,&-1,&0\end{smallmatrix}\right)\), and the latter into the former.

\Needspace{8\baselineskip}
\Article{276}
\textsc{Theorem.} \emph{The number of classes into which all ternary forms of a given determinant are divided is always finite.}

\SourceLeaf{319}
\emph{Proof.} I. The number of all forms \(\left(\begin{smallmatrix}a,&a',&a''\\b,&b',&b''\end{smallmatrix}\right)\) of a given determinant \(D\) in which \(a=0\), \(b''=0\), \(b\) is no greater than half the greatest common divisor of the numbers \(a'\), \(b'\); and \(a''\) no greater than \(b'\), is clearly finite. For since we must have \(a'b'b'=D\), the only possible values of \(b'\) are \(+1\), \(-1\) and the roots, with positive and negative signs, of the squares dividing \(D\) (if there are any others besides \(1\)). There are finitely many such values. For each value of \(b'\), however, the value of \(a'\) is determined, and the values of \(b\), \(a''\) are clearly limited to finitely many possibilities.

II. Similarly, there are finitely many forms \(\left(\begin{smallmatrix}a,&a',&a''\\b,&b',&b''\end{smallmatrix}\right)\) of determinant \(D\) in which \(a\) is not \(=0\) and is no greater than \(\frac{4}{3}\sqrt[3]{\pm D}\); \(b''b''-aa'=A''\) is not \(=0\) and is no greater than \(\frac{4}{3}\sqrt[3]{D^2}\); \(b''\) is no greater than \(\frac{1}{2}a\); and \(ab-b'b''=B\) and \(a'b'-bb''=B'\) are no greater than \(\frac{1}{2}A''\). For the number of all combinations of values of \(a\), \(b''\), \(A''\), \(B\), \(B'\) will be finite. Once these are individually determined, the remaining coefficients of the form, \(a'\), \(b\), \(b'\), \(a''\), and the coefficients of its adjoint
\[
bb-a'a''=A,\quad b'b'-aa''=A',\quad a''b''-bb'=B''
\]
will also be determined by the following equations:
\[
\begin{aligned}
a'&=\frac{b''b''-A''}{a},\quad A'=\frac{BB-aD}{A''},\\
A&=\frac{B'B'-a'D}{A''},\quad B''=\frac{BB'+b''D}{A''}\\
b&=\frac{AB-B'B''}{D}=-\frac{Ba'+B'b''}{A''},\\
b'&=\frac{A'B'-BB''}{D}=-\frac{Bb''+B'a}{A''}\\
a''&=\frac{b'b'-A'}{a}=\frac{bb-A}{a'}=\frac{bb'+B''}{b''}
\end{aligned}
\]
Now all those forms are obtained by choosing, from all combinations of values of \(a\), \(b''\), \(A''\), \(B\), \(B'\), those for which \(a'\), \(a''\), \(b\), \(b'\) also take integer values, so their number is clearly finite.

III. All the forms in I and II therefore constitute finitely many classes. There may even be fewer classes than forms if some of the forms are equivalent to one another. Now by the preceding investigations every ternary form of determinant \(D\) must be equivalent to one of those forms, \emph{that is}, belong to one of the classes which they constitute: these classes will therefore contain all forms of determinant \(D\), \emph{that is}, all ternary forms of determinant \(D\) will be divided into finitely many classes. \emph{Q.E.D.}
